\documentclass[12pt,a4paper]{article}
\usepackage{resumo}

\EnableReviewColors % Comente esta linha para remover as cores

\title{TÍTULO DESCRITIVO EM MAIÚSCULO}

\author{
    \reviewcolor{red}{Preencher após aprovação},
    Nome Autor,
    Coautor X,
    \dots,
    Orientador Fulano
}

\emails{
    \reviewcolor{red}{Preencher após aprovação},
    {email1, email2}@alu.ufc.br,
    outro@ufc.br
}

\begin{document}

\maketitle

\reviewcolor{purple}{
\textbf{Leia antes de escrever}: este modelo pode ser utilizado para trabalhos de pesquisa e relatos de experiência. O resumo deve descrever um trabalho realizado, apresentar resultados obtidos e ser escrito em um único parágrafo, com 2.000 a 3.000 caracteres, incluindo palavras-chave e ODS. A formatação é aplicada automaticamente pelo \LaTeX. Antes de gerar a versão final, desative as cores comentando \texttt{\textbackslash EnableReviewColors}.
}
\reviewcolor{purple}{
\textbf{Uso de IA}: ferramentas de IA generativa podem ser utilizadas como apoio à organização de ideias, revisão,
tradução e aprimoramento da redação. Os autores permanecem responsáveis pelo conteúdo e devem verificar as informações, dados, referências e afirmações produzidas ou sugeridas por essas ferramentas.
}
% Linhas em branco entre blocos geram um novo parágrafo; remova-as antes de realizar as submissões.

\reviewcolor{blue}{
\textbf{Contexto e Objetivo}: apresente brevemente o tema e o problema ou motivação do trabalho. Em seguida, indique claramente o objetivo geral, preferencialmente utilizando a palavra \textbf{objetivo}.
}
%
\reviewcolor{green}{
\textbf{Desenvolvimento}: explique como o trabalho foi realizado. Em pesquisas, apresente os principais procedimentos, participantes, materiais, técnicas e métodos de análise. Em relatos de experiência, descreva as atividades realizadas, o público envolvido, o período de execução e as estratégias adotadas. Inclua somente as informações necessárias para compreender como os resultados foram obtidos. Não confunda atividade com resultado: ``doze encontros foram realizados'' é um indicador de atividade, enquanto ``a taxa de retenção foi de 50\%'' é um indicador de resultado.
}
%
\reviewcolor{orange}{
\textbf{Resultados e Discussão}: apresente os principais resultados obtidos. Em trabalhos ainda em andamento, apresente os resultados parciais já obtidos. Não descreva apenas as atividades realizadas: evidencie seus resultados e, sempre que possível, utilize indicadores quantitativos e evidências qualitativas. Discuta brevemente a relação dos resultados com o objetivo proposto.
}
\reviewcolor{brown}{
\textbf{Conclusões e Encaminhamentos}: retome o objetivo e sintetize as principais conclusões, destacando as contribuições do trabalho. Quando pertinente, indique limitações, possibilidades de aplicação dos resultados ou perspectivas para trabalhos futuros.
}
\reviewcolor{black}{
\textbf{Exemplo}: este relato de experiência apresenta as atividades desenvolvidas no âmbito da monitoria da disciplina de Matemática, com o objetivo de apoiar os estudantes na consolidação dos conteúdos e reduzir as dificuldades observadas ao longo do semestre. As ações incluíram atendimentos presenciais e remotos, plantões semanais de dúvidas, resolução comentada de exercícios, elaboração de listas complementares e produção de materiais de apoio. Ao longo do período letivo, foram realizados 24 encontros de monitoria, com atendimento direto a 86 estudantes e participação média de 18 discentes por encontro. Para avaliar a efetividade das ações, foi aplicado um formulário eletrônico ao final do semestre, respondido por 61 participantes, contemplando questões sobre frequência, percepção de aprendizagem, clareza das explicações e contribuição da monitoria para o desempenho na disciplina. Os resultados indicaram que 93\% dos respondentes consideraram que a monitoria contribuiu para a compreensão dos conteúdos e 89\% afirmaram sentir maior segurança para resolver exercícios após participar das atividades. Além disso, os comentários qualitativos destacaram a disponibilidade do monitor, a organização dos materiais e a resolução passo a passo dos problemas como os principais pontos positivos. Conclui-se que as ações desenvolvidas atingiram o objetivo proposto, contribuindo para o processo de aprendizagem e indicando que a manutenção e ampliação da monitoria podem favorecer o desempenho dos estudantes em disciplinas com elevado índice de retenção.
\textbf{Atenção}: o exemplo serve apenas para demonstrar a estrutura e o nível de detalhamento esperados. Não copie seu conteúdo ou redação.
}

\vspace{0.5cm}

\textbf{Palavras-chave:}
Informe de três a cinco palavras-chave ou expressões curtas, iniciadas por letra minúscula e separadas por \textbf{ponto e vírgula}.

\textbf{ODS:}
Informe apenas um Objetivo de Desenvolvimento Sustentável (ODS), utilizando a
nomenclatura oficial em letra minúscula. Consulte:
\url{https://brasil.un.org/pt-br/sdgs}.

\end{document}

